
\documentclass{article}
\usepackage{colortbl}
\usepackage{makecell}
\usepackage{multirow}
\usepackage{supertabular}

\begin{document}

\newcounter{utterance}

\centering \large Interaction Transcript for game `imagegame', experiment `random\_grids', episode 6 with reka{-}flash{-}3.1{-}t0.0.
\vspace{24pt}

{ \footnotesize  \setcounter{utterance}{1}
\setlength{\tabcolsep}{0pt}
\begin{supertabular}{c@{$\;$}|p{.15\linewidth}@{}p{.15\linewidth}p{.15\linewidth}p{.15\linewidth}p{.15\linewidth}p{.15\linewidth}}
    \# & $\;$A & \multicolumn{4}{c}{Game Master} & $\;\:$B\\
    \hline

    \theutterance \stepcounter{utterance}  
    & & \multicolumn{4}{p{0.6\linewidth}}{
        \cellcolor[rgb]{0.9,0.9,0.9}{
            \makecell[{{p{\linewidth}}}]{
                \texttt{\tiny{[P1$\langle$GM]}}
                \texttt{Let us do the following task. The goal is to fill an empty grid that looks like this:} \\
\\ 
\texttt{▢ ▢ ▢ ▢ ▢} \\
\texttt{▢ ▢ ▢ ▢ ▢} \\
\texttt{▢ ▢ ▢ ▢ ▢} \\
\texttt{▢ ▢ ▢ ▢ ▢} \\
\texttt{▢ ▢ ▢ ▢ ▢} \\
\\ 
\texttt{A filled grid below is 5 by 5 and can look like this:} \\
\\ 
\texttt{▢ ▢ ▢ ▢ ▢} \\
\texttt{▢ ▢ E ▢ ▢} \\
\texttt{▢ ▢ ▢ ▢ ▢} \\
\texttt{▢ ▢ ▢ ▢ ▢} \\
\texttt{X X X X X} \\
\\ 
\texttt{I want you to describe this grid to me, step by step. You don't need to describe the empty squares, which are denoted with "▢". Only describe the location of letters in the grid. Then you wait for me to say "What is your next command?", and then you continue with the next step. Take the size of the grid into consideration while giving commands. You can give only one command at each turn but it can refer to single or multiple cells, rows, columns or even the whole grid at once. When you have described everything, you say "Command: DONE". Only give the command and omit everything else.} \\
\\ 
\texttt{For the filled grid above, here are the example steps.} \\
\\ 
\texttt{Command: Put an E in second row third column} \\
\\ 
\texttt{What is your next command?} \\
\texttt{Command: Fill the last row with X} \\
\\ 
\texttt{What is your next command?} \\
\texttt{Command: DONE} \\
\\ 
\texttt{Ok. Please do this for the following example, which is a 5 by 5 grid.} \\
\texttt{M ▢ ▢ ▢ ▢} \\
\texttt{M ▢ ▢ M ▢} \\
\texttt{▢ ▢ ▢ ▢ ▢} \\
\texttt{▢ ▢ M ▢ ▢} \\
\texttt{▢ M ▢ ▢ ▢} \\
\texttt{What is your next command?} \\
            }
        }
    }
    & \\ \\

    \theutterance \stepcounter{utterance}  
    & \multicolumn{4}{p{0.6\linewidth}}{
        \cellcolor[rgb]{0.9,0.9,0.9}{
            \makecell[{{p{\linewidth}}}]{
                \texttt{\tiny{[P1$\rangle$GM]}}
                \texttt{<reasoning>} \\
\texttt{Let's see. The user provided a 5x5 grid that's partially filled with Ms and some empty squares. They want me to describe the steps to fill this grid by specifying commands, similar to the example they gave. The example had commands like placing an E in a specific cell and filling a whole row with Xs. So, I need to figure out the positions of the Ms in the given grid and translate that into the appropriate commands.} \\
\\ 
\texttt{First, let me visualize the grid they provided:} \\
\\ 
\texttt{Row 1: M ▢ ▢ ▢ ▢} \\
\texttt{Row 2: M ▢ ▢ M ▢} \\
\texttt{Row 3: ▢ ▢ ▢ ▢ ▢} \\
\texttt{Row 4: ▢ ▢ M ▢ ▢} \\
\texttt{Row 5: ▢ M ▢ ▢ ▢} \\
\\ 
\texttt{So, the Ms are located at:} \\
\\ 
\texttt{{-} Row 1, Column 1} \\
\texttt{{-} Row 2, Column 1} \\
\texttt{{-} Row 2, Column 4} \\
\texttt{{-} Row 4, Column 3} \\
\texttt{{-} Row 5, Column 2} \\
\\ 
\texttt{Wait, let me check each row again to make sure I didn't miscount.} \\
\\ 
\texttt{Row 1: The first cell is M, the rest are empty. So (1,1) is M.} \\
\\ 
\texttt{Row 2: First cell M, then three empty, then M in column 4. So (2,1) and (2,4) are M.} \\
\\ 
\texttt{Row 3: All empty.} \\
\\ 
\texttt{Row 4: Third cell is M. So (4,3).} \\
\\ 
\texttt{Row 5: Second cell is M. So (5,2).} \\
\\ 
\texttt{So total of 5 Ms. The task is to describe the steps to place these Ms in the grid. The example used commands like "Put an E in second row third column" and "Fill the last row with X". So each command can target a single cell, multiple cells in a row/column, or the entire grid.} \\
\\ 
\texttt{In this case, since the Ms are scattered, maybe the best approach is to place each M individually. Alternatively, check if any of them can be grouped. Let's see:} \\
\\ 
\texttt{Looking at the positions:} \\
\\ 
\texttt{(1,1), (2,1), (2,4), (4,3), (5,2). These don't form a contiguous block or a row/column. So each would need individual commands. But maybe some can be addressed in a} \\
            }
        }
    }
    & & \\ \\

    \theutterance \stepcounter{utterance}  
    & & & \multicolumn{2}{p{0.3\linewidth}}{
        \cellcolor[rgb]{0.9,0.9,0.9}{
            \makecell[{{p{\linewidth}}}]{
                \texttt{\tiny{[GM$|$GM]}}
                \texttt{Invalid instruction format} \\
            }
        }
    }
    & & \\ \\

\end{supertabular}
}

\end{document}
